% Part: propositional-logic
% Chapter: propositional-logic
% Section: preliminaries

\documentclass[../../../include/open-logic-section]{subfiles}

\begin{document}

\olfileid[pt]{pl}{syn}{pre}

\olsection{Preliminares}

\begin{thm}[\emph{Princípio da indução sobre !!{formula}s}]
  \ollabel{thm:induction}  
  Se alguma propriedade~$P$ vale para todas as !!{formula}s atômicas e é
  tal que
  \begin{enumerate}
    \tagitem{prvNot}{ela vale para $\lnot !A$ sempre que vale
      para~$!A$;}{}
    \tagitem{prvAnd}{ela vale para $(!A \land !B)$
      sempre que vale para $!A$ e~$!B$;}{}
    \tagitem{prvOr}{ela vale para $(!A \lor !B)$
      sempre que vale para $!A$ e~$!B$;}{}
    \tagitem{prvIf}{ela vale para $(!A \lif !B)$
      sempre que vale para $!A$ e~$!B$;}{}
    \tagitem{prvIff}{ela vale para $(!A \liff !B)$
      sempre que vale para $!A$ e~$!B$;}{}
  \end{enumerate}
  então $P$ vale para todas as !!{formula}s.
\end{thm}

\begin{proof}
  Seja $S$ a coleção de todas as !!{formula}s com a
  propriedade~$P$. Claramente, $S \subseteq \Frm[L_0]$. $S$~satisfaz todas as
  condições de \olref[fml]{defn:formulas}: contém todas as
  !!{formula}s atômicas e é fechada sob os !!{operator}s. $\Frm[L_0]$ é
  a menor classe desse tipo, portanto $\Frm[L_0] \subseteq S$. Logo, $\Frm[L_0] = S$, e
  toda fórmula tem a propriedade~$P$.
\end{proof}

\begin{prop}\ollabel{prop:balanced}
   Toda !!{formula} em~$\Frm[L_0]$ é \emph{balanceada}, pois tem
   tantos parênteses esquerdos quanto direitos.
\end{prop}

\begin{prob}
Demonstre \olref[pl][syn][pre]{prop:balanced}
\end{prob}

\begin{prop} \ollabel{prop:noinit}
Nenhum segmento inicial próprio de !!a{formula} é !!a{formula}.
\end{prop}

\begin{prob}
  Demonstre \olref[pl][syn][pre]{prop:noinit}
\end{prob}

\begin{prop}[Leitura Unívoca]
Toda !!{formula}~$!A$ em $ \Frm[L_0]$ tem exatamente uma análise sintática em uma das
seguintes formas:
\begin{enumerate}
\tagitem{prvFalse}{$\lfalse$.}{}

\tagitem{prvTrue}{$\ltrue$.}{}

\item $\Obj p_n$ para algum $\Obj p_n \in  \PVar$.
  
\tagitem{prvNot}{$\lnot !B$ para alguma !!{formula}~$!B$.}{}

\tagitem{prvAnd}{$(!B \land !C)$ para algumas !!{formula}s $!B$ e~$!C$.}{}

\tagitem{prvOr}{$(!B \lor !C)$ para algumas !!{formula}s $!B$ e~$!C$.}{}

\tagitem{prvIf}{$(!B \lif !C)$ para algumas !!{formula}s $!B$ e~$!C$.}{}

\tagitem{prvIff}{$(!B \liff !C)$ para algumas !!{formula}s $!B$ e~$!C$.}{}
\end{enumerate}
Além disso, essa análise sintática é \emph{única}.
\end{prop}

\begin{proof}
Por indução sobre $!A$. Por exemplo, suponha que $!A$ tenha duas
leituras distintas como $(!B \lif !C)$ e $(!B' \lif !C')$. Então, $!B$ e $!B'$
devem ser iguais (caso contrário, um seria um segmento inicial próprio do
outro); assim, se as duas leituras de $!A$ forem distintas, isso deve ocorrer porque
$!C$ e $!C'$ são leituras distintas da mesma sequência de símbolos,
o que é impossível pela hipótese de indução.
\end{proof}


\begin{defn}[Substituição Uniforme]
Se $!A$ e $!B$ são !!{formula}s e $\Obj p_i$ é uma !!{propositional
variable}, então $\Subst{!A}{!B}{\Obj p_i}$ denota o resultado de
substituir cada ocorrência de $\Obj p_i$ por uma ocorrência de $!B$ em $!A$;
de modo semelhante, a substituição simultânea de $\Obj p_1$, \dots,~$\Obj p_n$ por
!!{formula}s $!B_1$, \dots,~$!B_n$ é denotada por
$\SSubst{!A}{\subst{!B_1}{\Obj p_1},\dots,\subst{!B_n}{\Obj p_n}}$.
\end{defn}

\begin{prob} Para cada uma das cinco !!{formula}s abaixo, determine se a 
 !!{formula} pode ser expressa como uma substituição \( \Subst{!A}{!B}{\Obj p_i} \),
 em que \( !A \) é (i) \( \Obj p_0 \); (ii) \( ( \lnot \Obj p_0 \land \Obj
 p_1) \); e (iii) \( ( ( \lnot \Obj p_0 \lif \Obj p_1 ) \land \Obj
 p_2 ) \). Em cada caso, especifique a substituição pertinente.
  \begin{enumerate}
    \item \( \Obj p_1 \) 
    \item \( ( \lnot \Obj p_0 \land \Obj p_0 ) \)
    \item \( ( ( \Obj p_0 \lor \Obj p_1 ) \land \Obj p_2 )  \)
    \item \( \lnot ( ( \Obj p_0 \lif \Obj p_1 ) \land \Obj p_2 ) \)
    \item \( (( \lnot ( \Obj p_0 \lif \Obj p_1 ) \lif ( \Obj p_0 \lor \Obj p_1 )) \land \lnot ( \Obj p_0 \land \Obj p_1 )) \)
  \end{enumerate}
\end{prob}

\begin{prob}
Dê uma definição matematicamente rigorosa de $\Subst{!A}{!B}{p}$ por
indução.
\end{prob}

\end{document}

